\subsection{Sources and available records}
\label{app:reproducibility}
\label{sec:b-artifacts}
The \href{https://github.com/lrlorry/8800_report/blob/001579c9771759b672ef377b10dbc3dd959c7979/docs/artifact-access.md}{project archive}
provides source, settings, timings and selected inputs. Complete vectors,
trained states and historical executable hashes are not all available.
Deleted candidate arrays from the strict selection check require regeneration for an independent
re-sort. \Cref{tab:a-components} maps versions to component controls;
their speedups cannot be multiplied into a whole-system result.
Sampling and generalisation limits follow \cref{sec:6-2,sec:6-3,sec:8-3}.

\Cref{tab:3-campaigns} gives the names used in later chapters; record labels locate the raw data in the project archive.

\par\medskip\noindent\begin{minipage}{\linewidth}
\centering\captionsetup{hypcap=false}
\captionof{table}[Experiments and their records]{Experiments and their records.}
\label{tab:3-campaigns}
\footnotesize\setlength{\tabcolsep}{3pt}
\begin{tabularx}{\linewidth}{@{}>{\raggedright\arraybackslash}p{.24\linewidth}l>{\raggedright\arraybackslash}Xl@{}}
\toprule
Name in the text & Record & What it measures & Results \\
\midrule
Main comparison & E1 & JHQ-GPU and five baselines; four datasets, batch 1,000 & \cref{sec:7-2} \\
Precision check & D & FP32 versus TF32; archived baseline measurements & \cref{sec:7-precision} \\
Same-session check & T1 & JHQ and baselines measured together & \cref{sec:b-precision-anchors} \\
Strict selection check & T4 & CPU re-sort of GPU-computed candidate scores & \cref{sec:7-strict-selection} \\
Small-batch study & T5 & Cross-system search at batches 8 and 32 & \cref{sec:7-smallbatch} \\
Repeated cold builds & A1 & Build time and sampled memory peaks & \cref{sec:7-7} \\
Instruction check & A2 & Compiled shared-bound read/update instructions & \cref{sec:5-5} \\
Alignment check & A3 & Row alignment of the two OpenAI exports & \cref{sec:7-2} \\
Budget decision traces & B & Calibration choices and complete enumeration & \cref{sec:7-budget-trace} \\
Representation audit & C1 & Original-input recoding and placement checks & \cref{sec:4-representation-audit} \\
Centring test & T2 & Original versus mean-centred representations & \cref{sec:8-centring} \\
\bottomrule
\end{tabularx}
\end{minipage}\par\medskip

\par\medskip\noindent\begin{minipage}{\linewidth}
\centering
\captionsetup{hypcap=false}
\captionof{table}[Component implementations]{Implementation families used by component tests. Each experiment retains its own control; versions are not successive whole-system speedups.}
\label{tab:a-components}
\small\setlength{\tabcolsep}{4pt}
\begin{tabularx}{\linewidth}{@{}l>{\raggedright\arraybackslash}X@{}}
\toprule
Version & Mechanism / experimental role\\\midrule
v57 & Earlier GPU implementation in the historical CPU comparison\\
v61 & Word-load and lane-refinement controls\\
v62 / v63 & Grouped full-key selection / shared bound\\
v65 / v69 & Batch/skew scheduling / radix probe selection\\
v69 / v71 / v73 & Budget calibration / LUT sweep / strict-score export\\
v69 / v72 & OpenAI: TF32 versus FP32 add/query arithmetic on the same trained cache\\
\bottomrule
\end{tabularx}
\end{minipage}\par\medskip

\subsection{Timing and calibration conventions}
\label{sec:a-e1-protocol}
The evaluation explicitly enables sixteen-lane refinement and automatic
scan scheduling among one block and 2/4/8 groups, with shared bounds for
multiple groups. The library defaults to scalar refinement and one block;
the optional quota-based approximate selector is disabled. One-time
scheduler calibration is excluded from steady-state timing. These scheduler-only ranges cover
pseudoquery gathering, the first calibration call and the median of repeated calls:
\[
 T_{\rm gather}\in[0.81,7.04]\,\mathrm{ms},\quad
 T_{\rm first}\in[6.00,27.77]\,\mathrm{ms},\quad
 \operatorname{median}(T_{\rm repeat})\in[2.72,18.13]\,\mathrm{ms}.
\]

For the main comparison, JHQ warms once per variant/round and baselines warm three times
initially. Driver just-in-time (JIT) compilation and JHQ CUDA Graph
capture, which records GPU operations for replay, precede timing.
Slow baseline configurations use fewer timed calls based on a warm-up
estimate. Repeat-call ID-order changes can include tied reordering and
exclude no points by themselves; this differs from duplicate IDs within
one result. JHQ lacks the same repeat-call diagnostic flag.

The packed-load comparison retains one aggregated timing per variant.
Sampled-budget tables combine the median selected budget with mean quality
over 64 samples; mean recall is reconstructed from printed best recall
minus logged mean loss. These summaries need not describe one fixed budget.

\label{sec:a-budget-contract}
\label{sec:a-supplement}
\label{sec:a-t4-t5}
Budget experiments use one warm-up and three timed calls per fixed
budget. Calibration time includes the reference search, searches at trial
budgets and the host-side comparison of their results ($O(Sk^2)$). It
excludes gathering the sample queries, the one-time scheduler calibration
and graph capture. Warm-up runs over all queries, but its outputs are
not used for budget selection, and calibration never uses GT. Selected-budget QPS uses fixed-budget timing; direct
reuse tests follow \cref{sec:6-5}. For 64 sorted recall losses, the
reported percentile uses
\[
 L_{95}=L^{\rm sorted}_{\lfloor0.95(64-1)\rfloor}
       =L^{\rm sorted}_{59}\qquad\text{(zero-based)}.
\]
